\chapter{Evaluation}

The evaluation is organized around claims. Each graph answers one question and
is interpreted immediately.

All numeric claims in this chapter come from the frozen data snapshot in
\code{data/snapshot/}. The snapshot records the exact CSV and Markdown files
used to generate the current graph set. This avoids quoting numbers from
mutable \code{latest} symlinks.

\section{Methodology}

All backends run under the same outer benchmark harness. The workload is
TPC-C. The sweep varies Mako worker thread count from \(t=1\) to \(t=11\). The
main no-disk comparison includes Single-Raft without ReplayPool, Single-Raft
with 1:1 ReplayPool, Multi-Raft, and Paxos.

The headline metric is \code{throughput\_ops\_sec}, extracted from the leader-side
benchmark process. Worker CPU metrics show whether transaction workers are
active. Replay counters show follower progress. The simulated-disk counters are used
only for persistence experiments.

\section{No-Disk Scalability}

This experiment asks whether Raft can stay in Paxos's performance class when
disk persistence is disabled. Multi-Raft tests the conservative hypothesis:
Raft should match Paxos when given a similar partition-level topology.
Single-Raft tests the consolidation hypothesis: one consolidated Raft log,
with one local replica running in each process, can be viable if the replay
path is engineered correctly.

\begin{figure}[htbp]
\centering
\includegraphics[width=0.96\linewidth]{figures/graphs/fig07_no_disk_scalability.pdf}
\caption[No-disk scalability]{No-disk scalability. The left panel reports leader-side committed throughput; the right panel reports follower replay throughput.}
\label{fig:no-disk-scalability}
\end{figure}

Figure~\ref{fig:no-disk-scalability} separates leader-side committed
throughput from follower replay throughput. At \(t=11\), Single-Raft without
ReplayPool reports about 593K transactions/s, Single-Raft with ReplayPool
reports about 526K transactions/s, Multi-Raft reports about
510K transactions/s, and Paxos reports about 472K transactions/s. A throughput-only reading would make the no-ReplayPool
configuration look best.

The right panel shows why that reading is wrong. At \(t=11\), Single-Raft
without ReplayPool reaches only about 143K transactions/s of follower replay
progress. Single-Raft with ReplayPool reaches about 499K transactions/s;
Multi-Raft reaches about 484K transactions/s; Paxos reaches about
450K transactions/s. ReplayPool therefore makes
Single-Raft much closer to a fully replayed replicated path. The result is not
that ReplayPool maximizes leader-side throughput. The result is that
ReplayPool makes the Single-Raft measurement credible as replicated execution.

This distinction is central to the thesis. Leader-side throughput is the
headline benchmark metric, but follower replay progress determines whether
the replicated system is actually keeping up with the committed history.

\section{Worker CPU}

The next question is whether the benchmark workers are doing real work or
whether the replication path causes the system to idle. Worker CPU does not
prove correctness by itself, but it helps interpret throughput results.

\begin{figure}[htbp]
\centering
\includegraphics[width=0.96\linewidth]{figures/graphs/fig08_worker_cpu_utilization.pdf}
\caption[Worker CPU utilization]{Worker CPU utilization across the no-disk scalability sweep.}
\label{fig:worker-cpu-utilization}
\end{figure}

Figure~\ref{fig:worker-cpu-utilization} shows that worker CPU remains above
one full core per active worker through the high-thread part of the sweep. At
\(t=11\), the mean worker CPU is about 166\% for Single-Raft with ReplayPool,
253\% for Multi-Raft, and 222\% for Paxos. These numbers support the interpretation that the benchmark
is exercising Mako's transaction workers rather than measuring an idle
replication harness.

The CPU curves also explain why throughput comparisons should not be reduced
to one line. Multi-Raft and Paxos spend more CPU per active worker at high
thread counts, while Single-Raft with ReplayPool keeps the consolidated Raft
path competitive. The thesis claim is not that one backend is universally
faster. The claim is that Raft can remain in the same performance class when
the pipeline is preserved.

\section{ReplayPool Mechanism}

This experiment isolates the ReplayPool mechanism by varying both Mako worker
threads and \code{MAKO\_REPLAY\_THREADS}. The replay thread counts are 0, 1, 2,
4, 8, and 11, and each configuration is swept from \(t=1\) to \(t=11\).
The question is whether replay parallelism changes useful follower progress,
not whether more replay threads monotonically increase leader-side throughput.

\begin{figure}[htbp]
\centering
\includegraphics[width=0.96\linewidth]{figures/graphs/fig09_replay_pool_matrix.pdf}
\caption[ReplayPool sensitivity]{ReplayPool sensitivity across worker counts and replay-pool sizes.}
\label{fig:replaypool-matrix}
\end{figure}

Figure~\ref{fig:replaypool-matrix} shows that leader-side throughput stays in a
relatively narrow band at high thread counts, even when the replay pool is too small. At \(t=11\),
throughput ranges from about 563K transactions/s to 584K transactions/s across
the tested replay-pool sizes. A leader-throughput-only graph would therefore
miss the actual mechanism.

The replay panels show the missing evidence. At \(t=11\), zero replay threads
produce about 137K transactions/s of follower replay progress, or
about 23.5\% of the leader-side throughput. One replay thread is similar at
about 133K transactions/s and 22.8\% completeness. Two replay
threads improve follower progress to about 257K transactions/s, but
that is still only about 44.6\% completeness. Four replay threads reach
about 461K transactions/s and 81.5\% completeness. Eight and eleven
replay threads reach about 541K transactions/s and 534K transactions/s,
respectively, and both stay near 95\% completeness.

The same pattern appears before the largest endpoint. At \(t=6\), zero and one
replay thread keep only about 40\% of the leader's committed history replayed
on the follower. Two replay threads reach about 79\%. Four, eight, and eleven
replay threads are all near 95\%. The threshold therefore grows with worker
parallelism: a small replay pool may look adequate at low thread counts, but
falls behind as the leader submits more committed work.

This is the reason the thesis uses the 1:1 ReplayPool configuration for the
main Single-Raft comparison. Here, 1:1 means \code{MAKO\_REPLAY\_THREADS} equals the Mako
worker-thread count. It is a conservative setting that keeps follower replay
close to leader-side commit throughput across the tested sweep. ReplayPool is
therefore best understood as a replay-completeness mechanism, not as a
monotonic leader-throughput optimization.

\begin{figure}[htbp]
\centering
\includegraphics[width=0.96\linewidth]{figures/graphs/fig10_replay_pool_matrix_heatmap.pdf}
\caption[Follower replay completeness]{Follower replay completeness across worker counts and replay-pool sizes.}
\label{fig:replaypool-heatmap}
\end{figure}

Figure~\ref{fig:replaypool-heatmap} is a compact way to read the same result. The dark low-completeness
region is concentrated where worker count is high and replay-pool size is too
small. The high-completeness region appears when the replay pool is large
enough to track the leader's committed history. This figure is the mechanism
evidence behind the statement that Single-Raft is viable only after replay
parallelism is engineered explicitly.

\section{Simulated Persistence Throughput}

This experiment asks how the backends behave when persistence cost is modeled.
The three conditions are no disk, an NVMe-like simulated-disk model, and a
Cloud-SSD-like simulated-disk model.

\begin{figure}[htbp]
\centering
\includegraphics[width=0.96\linewidth]{figures/graphs/fig11_disk_persistence_throughput.pdf}
\caption[Simulated persistence throughput]{Simulated persistence throughput under no-disk, NVMe-like, and Cloud-SSD-like settings.}
\label{fig:disk-persistence-throughput}
\end{figure}

Figure~\ref{fig:disk-persistence-throughput} shows that the NVMe model stays
close to the no-disk curves. At \(t=11\), Single-Raft with
ReplayPool reaches about 545K transactions/s under NVMe, Multi-Raft about
509K transactions/s, and Paxos about 475K transactions/s. Those values are close to the no-disk
ordering and support the interpretation that this disk model does not dominate
the high-thread runs.

The Cloud-SSD model behaves differently. At high thread counts, all three
backends converge near 340K transactions/s. At \(t=11\), Single-Raft
with ReplayPool reports about 340K transactions/s, Multi-Raft about
341K transactions/s, and Paxos about 341K transactions/s. This convergence is important: it
does not distinguish the consensus backends. It suggests that, under the
Cloud-SSD simulated-disk model, the storage path becomes the shared limiting factor.

This is a simulated persistence result. It should be read as evidence about
the modeled storage pressure, not as a claim about production SSD behavior.

\section{Simulated-Disk Proof}

The disk-throughput graph is only persuasive if it is paired with storage-work
evidence. The simulated-disk instrumentation provides that evidence by
recording the bytes and writes charged during the benchmark.

\begin{figure}[htbp]
\centering
\includegraphics[width=0.96\linewidth]{figures/graphs/fig12_simulated_disk_proof_cloudssd.pdf}
\caption[Cloud-SSD simulated-disk proof]{Cloud-SSD flattening with simulated-disk proof statistics.}
\label{fig:simulated-disk-proof-cloudssd}
\end{figure}

Figure~\ref{fig:simulated-disk-proof-cloudssd} shows that, at high thread
counts, the Cloud-SSD throughput curves flatten near 340K transactions/s while
the simulated-disk counters report real measured storage work. At \(t=11\),
all three backends write about 7.49 GB through the simulated disk and perform
about 24.6K writes. The normalized value is about 733 simulated-disk
bytes per committed transaction.

\begin{figure}[htbp]
\centering
\includegraphics[width=0.96\linewidth]{figures/graphs/fig13_cloudssd_bytes_per_txn.pdf}
\caption[Cloud-SSD bytes per transaction]{Cloud-SSD simulated-disk bytes per committed transaction.}
\label{fig:cloudssd-bytes-per-txn}
\end{figure}

Figure~\ref{fig:cloudssd-bytes-per-txn} shows that the storage work per
committed transaction is stable across backends and rises modestly at high
thread counts. This supports the explanation that the Cloud-SSD flattening is
consistent with measured persistence work. It does not prove real SSD
behavior; it proves that the controlled simulated-disk model was doing comparable
measured work across the three backends when the throughput curves converged.

\section{Evaluation Summary}

The evaluation supports four conclusions.

First, Raft can replace Paxos in Mako without falling out of Paxos's
performance class. Multi-Raft is the direct Paxos-shaped replacement and
tracks the original backend closely.

Second, Single-Raft can be competitive, but only if interpreted with follower
replay progress. Without ReplayPool, high leader-side throughput hides the
fact that followers replay much more slowly. With ReplayPool, follower replay
progress approaches leader-side throughput.

Third, worker CPU and replay CPU counters make the mechanism visible. The
benchmark workers are active, and ReplayPool moves replay work out of the
apply path.

Fourth, disk claims require simulated-disk proof statistics. The Cloud-SSD model
flattens near 340K transactions/s, and that flattening is accompanied
by measured bytes, writes, and bytes per committed transaction. The thesis can
therefore explain the disk result without pretending that simulated disk is a full
production storage evaluation.

Three no-disk \(t=11\) reruns had coefficient of variation below 0.6\% for
Single-Raft, Multi-Raft, and Paxos. This repeatability result is supporting
evidence, not a main thesis figure.
